% Introduction — SIREN
% CITE markers indicate where a reference must be inserted before submission.
\section{Introduction}

Descriptor-based property prediction remains the cheapest route from molecular
structure to a numerical endpoint. A modern 2D descriptor set such as the RDKit
family contains a few hundred columns, most of them correlated by construction:
partial charges, surface-area contributions and fragment counts are all
integrals of the same underlying graph. Reducing that matrix to a handful of
interpretable columns is therefore routine, both to fit small datasets and to
report which structural features drive a prediction.

The reduction step is usually a marginal scorecard. Every descriptor is ranked
by its univariate correlation with the endpoint, the strongest are kept, and
highly correlated pairs are pruned. The procedure is fast, deterministic and
widely implemented, but it answers a question that is not the one being asked.
It measures how much a descriptor knows \emph{on its own}; a model uses how much
a descriptor adds \emph{given the others}. The two answers differ in three
situations that are common in chemistry.

First, \emph{suppression}. A descriptor can be strongly associated with the
endpoint and still show a near-zero marginal correlation when a correlated
partner already carries the same variation. A phenol hydroxyl fragment, for
instance, is a direct determinant of aqueous solubility, yet its univariate
association with $\log S$ is weak because the polarity it encodes is already
summarised by a global lipophilicity descriptor. A marginal ranking discards it;
a conditional ranking keeps it.

Second, \emph{synergy}. Some structural features matter only in combination ---
a hydrogen-bond donor is informative about permeability when an acceptor is also
present, and uninformative otherwise. No marginal statistic can see this,
because the information exists only in the joint distribution. A greedy
selection that rewards interaction against an already selected partner can in
principle find such pairs, but only if one member is selected first, which the
marginal evidence cannot motivate.

Third, \emph{instability}. Univariate rankings on a few hundred molecules are
noisy. A descriptor that leads the ranking on the observed sample may not do so
on a resample, and a selector that ignores this variance will happily spend its
budget on an artefact. The problem is not resolvable by a larger significance
threshold, because the ranking is computed once and its variability is never
estimated.

Separately, the evaluation of selectors is frequently optimistic. Selecting
descriptors on the full dataset and then evaluating a model on a held-out split
lets the held-out rows influence the selection, which inflates the reported
score; the inflation is invisible in a single train/test protocol and difficult
to compare across papers because it depends on how the split was drawn. A
selector is only credible if it is refitted inside every training fold.

This paper makes three contributions. We describe SIREN, a greedy forward
selector whose objective is conditional on the selected context and combines
four terms: a conditional mutual-information relevance, an interaction term
built from an orthogonalised product basis, a resampling stability term and an
optional redundancy penalty. We add two mechanisms without which the objective
is not usable on real descriptor matrices: a near-duplicate screen that prevents
a copied column from consuming a slot, and a permutation-calibrated synergy seed
that makes pure interaction pairs reachable. And we evaluate the selector with a
protocol that separates four questions that are usually conflated: whether it
recovers a known structure, whether it improves a held-out readout, how much
optimism leakage would have bought, and whether its selection is reproducible.

We report what the experiments actually show, including the cases where SIREN does
not help: on purely marginal or pure-noise designs it matches the scorecard, and
on interaction designs its recovery is high but not perfect. The code and the
selection traces are released so that both the successes and the failures can be
reproduced.
